\chapter*{Preface}

\section*{Why this book}

Most introductions to K3 surfaces begin with the surface --- a quartic in
$\PP^3$, a Kummer surface, a double plane branched along a sextic --- and
arrive at the lattice $H^2(X,\Z)$ as a piece of topology one computes about
it. This book runs the other way. It begins with integral lattices, because
for K3 surfaces the lattice is not a shadow of the geometry: by the Torelli
theorem it \emph{is} the geometry, up to the position of a single line in a
fixed $22$-dimensional space. A reader who has internalized how a rank-3
lattice $\U\oplus\lat{2n}$ sits inside $\U^3\oplus E_8(-1)^2$, how a vector of
norm $-2$ in it splits it, and how the isometries of that rank-3 lattice
realize a modular group, already knows the skeleton of the theory of
one-parameter families of lattice-polarized K3 surfaces. The rest is filling
in the flesh --- Hodge theory, periods, fibrations, differential equations ---
and knowing which classical theorem carries each step.

\section*{Who it is for}

The reader we have in mind has taken linear algebra seriously (including
bilinear forms over $\Z$), knows what a Riemann surface and a holomorphic
form are, and has met the words ``divisor'', ``line bundle'' and
``cohomology'' without necessarily being at ease with them. Chapters~1 and~7
require only the first of these. Chapters~2--6 use the rest, but every
theorem is stated in full and every proof that we do not give is located
precisely in a standard reference. We have tried to write the book we wished
we had had when the mathematics of \S\ref{ch:frontier} was being assembled
piece by piece from those references.

\section*{The status convention}

Every mathematical statement in this book carries one of four marks.

\begin{itemize}[leftmargin=2em]
\item \Lstatus\ \emph{Literature.} A theorem of the published record, quoted
  with a precise citation. We do not reprove it; where a proof is short and
  illuminating we include it and say so.
\item \Estatus\ \emph{Exact computation.} A finite computation over $\Q$ or
  $\Z$ that the reader can rerun; typically a Gram matrix, a determinant, an
  order of vanishing, or a $q$-expansion checked to a stated order. We never
  write ``verified'' for a finite check without stating its order.
\item \Nstatus\ \emph{Numerical.} A high-precision floating-point computation
  followed by recognition of an exact value. This is evidence, not proof, and
  the book says so wherever it uses one.
\item \Kstatus\ \emph{Kernel-checked.} A statement that has been formalized
  in Lean~4 and accepted by its kernel, in a public repository, depending only
  on the three standard axioms of the library. Chapter~\ref{ch:formal}
  explains exactly what this mark certifies --- and its two well-documented
  ways of certifying less than a name suggests.
\end{itemize}

A theorem marked \Lstatus\ is not less true than one marked \Kstatus. The
marks record \emph{how we know}, and a reader who wants to build on a
statement should know which kind of knowing it is. The habit of writing the
mark down is the single most useful thing this book has to teach.

\section*{How to read it}

Chapter~\ref{ch:lattices} is the foundation and should be read first; its
worked examples are the lattices that return in every later chapter.
Chapters~\ref{ch:k3}--\ref{ch:torelli} are the classical theory of K3
surfaces, told so that the lattice does the work. Chapter~\ref{ch:elliptic}
is the one place where explicit equations appear, because elliptic fibrations
are how one \emph{computes} a N\'eron--Severi lattice by hand.
Chapters~\ref{ch:pf} and~\ref{ch:modular} bring in analysis --- Picard--Fuchs
equations and modular curves --- and Chapter~\ref{ch:singular} is where they
meet the lattice at the points of Picard number $20$. Chapter~\ref{ch:frontier}
assembles the frontier: what is known about Cooper's sporadic sequences and
the $M_7$-polarized K3 family attached to them, with every claim carrying its
mark. Chapter~\ref{ch:formal} is about the \Kstatus\ mark itself.

Exercises are placed where the material is; hints are in
Appendix~\ref{app:hints}. An exercise marked $\star$ is one whose answer the
book uses later.

\section*{What this book is not}

It is not a research monograph: Huybrechts' \emph{Lectures on K3
Surfaces}~\cite{Huybrechts2016} and Barth--Hulek--Peters--Van de Ven
\cite{BHPV2004} remain the references for proofs, and we lean on them
constantly. It is not a treatise on formal verification; Chapter~\ref{ch:formal}
is deliberately short. And it contains no physics: the lattice
$\U\oplus\lat{2n}$ appears in string theory as well, and a reader from that
direction will recognize some matrices, but every sentence here is a
statement about mathematics and nothing else.

\section*{Provenance}

This edition was drafted with the assistance of an AI model (Claude, Fable
5.1) working inside the repository that holds the Lean development of
Chapter~\ref{ch:formal}. The kernel-checked statements are transcribed from
that repository at a stated commit; the literature statements are the
author's responsibility and have been checked against the cited sources to the
extent stated in each chapter's closing note. Corrections are welcome and
will be recorded, not silently applied.

\bigskip
\noindent\textit{Edition 0.1, 29 September 2026.}
